% Section content template for individual Educative lessons
% This file contains the actual content for a specific section
% Generated from Educative API section content

% Section: Quiz: Object-oriented Basics
% Section ID: 5962355341262848
% Section Slug: quiz-object-oriented-basics
% Generated: 2025-12-16T12:49:54.815479

\subsection{Match the answers}\label{khiQy5hFp6HJMSGOolcSw}
\\
As a game developer, you create popular games such as Tic-Tac-Toe and Sudoku. You use object-oriented programming (OOP) principles to design a well-structured and scalable project. OOP helps you define relationships between different components and organize your code efficiently.
\\
To get started, you establish a base class, \texttt{Game}, that contains common attributes and methods. Then , you create specific game classes---\texttt{TicTacToe} and \texttt{Sudoku}---which inherit from the \texttt{Game} class and add their unique attributes and behaviors. An example of how this project's structure is defined is given below:

\begin{lstlisting}[language=plaintext]
-> Game : Base Class
    -> playerNum : int              (Stores the number of players)
    -> makeMove() : void            (Represents making a move in the game)
    -> getPlayers() : string        (Returns player details)

-> TicTacToe : Derived from Game
    -> isHuman : bool               (Indicates whether the player is human)
    -> getIsHuman() : bool          (Returns whether the player is human)

-> Sudoku : Derived from Game
    -> value : int                  (Represents a value in the Sudoku grid)
    -> setValue() : void            (Allows setting a value in the Sudoku grid)
\end{lstlisting}

Below are four core OOP principles. Your task is to match each principle with the correct example based on the given class structure.
\begin{quote}
\textbf{Note:} Select an OOP principle in the left-hand column by clicking it, and then click one of the options in the right-hand column.
\end{quote}

\textit{Component type 'MatchTheAnswers' not yet supported.}

\subsection{MCQs}\label{E99f25gGckVKDrZUsgKEr}
\\
Challenge yourself by solving the following quiz questions.

\subsection*{Object-Oriented Basics}
\vspace{0.3cm}
\noindent
\textbf{Question 1:} What best captures the main advantage of using object-oriented programming (OOP)?
\\[0.2cm]
\begin{itemize}
 \item It focuses on writing low-level, memory-optimized code.
 \item \textbf{[CORRECT]} It groups related data and behavior into objects to model real-world systems more naturally.
\\[0.1cm]
\textit{Explanation:} 
OOP is about modeling software using objects, which often represent real-world things and encapsulate both data and behavior.
 \item It allows you to avoid writing functions entirely.
 \item It replaces the need for planning and system design.
\end{itemize}
\vspace{0.3cm}
\noindent
\textbf{Question 2:} Which statement is \emph{true} about objects and classes in OOP?
\\[0.2cm]
\begin{itemize}
 \item A class is a specific instance of an object.
 \item Objects can exist without a class.
 \item \textbf{[CORRECT]} A class defines the blueprint for creating objects.
\\[0.1cm]
\textit{Explanation:} 
A class provides the template or blueprint for creating objects, which are specific instances of the class.
 \item An object defines the template for creating classes.
\end{itemize}
\vspace{0.3cm}
\noindent
\textbf{Question 3:} What is the main purpose of encapsulation in OOP?
\\[0.2cm]
\begin{itemize}
 \item \textbf{[CORRECT]} To hide internal data of a class and expose only what's necessary.
\\[0.1cm]
\textit{Explanation:} 
Encapsulation involves wrapping data (variables) and behavior (methods) into a single unit and hiding internal details from external access.
 \item To make the code run faster.
 \item To allow one class to inherit from another.
 \item To combine two classes into one.
\end{itemize}
\vspace{0.3cm}
\noindent
\textbf{Question 4:} Why is it recommended to use private variables in a class?
\\[0.2cm]
\begin{itemize}
 \item So other classes can modify them freely.
 \item To reduce memory usage.
 \item \textbf{[CORRECT]} To prevent accidental or unauthorized changes to the data.
\\[0.1cm]
\textit{Explanation:} 
Making variables private protects them from outside interference and maintains data integrity.
 \item To make the class load faster.
\end{itemize}
\vspace{0.3cm}
\noindent
\textbf{Question 5:} What best represents encapsulation?
\\[0.2cm]
\begin{itemize}
 \item \textbf{[CORRECT]} A phone's interface showing apps but hiding internal circuits.
\\[0.1cm]
\textit{Explanation:} 
Like a phone, encapsulation shows a usable interface while hiding the complex internal mechanisms.
 \item A chef showing every recipe ingredient to customers.
 \item A book where every chapter is public.
 \item A classroom where every student's grade is shared openly.
\end{itemize}
\vspace{0.3cm}
\noindent
\textbf{Question 6:} What does abstraction mean in the context of OOP?
\\[0.2cm]
\begin{itemize}
 \item Hiding the real-world objects.
 \item \textbf{[CORRECT]} Focusing on essential features while hiding complex details.
\\[0.1cm]
\textit{Explanation:} 
Abstraction allows a programmer to deal with ideas rather than specific events---simplifying complexity by exposing only what's necessary.
 \item Storing multiple values in one object.
 \item Inheriting from multiple classes.
\end{itemize}
\vspace{0.3cm}
\noindent
\textbf{Question 7:} You're building a ride booking app. As a user, you can simply call \texttt{bookRide()} without knowing how the matching algorithm or pricing logic works.
\\
What OOP principle does this demonstrate?
\\[0.2cm]
\begin{itemize}
 \item Inheritance
 \item Encapsulation
 \item \textbf{[CORRECT]} Abstraction
\\[0.1cm]
\textit{Explanation:} 
Abstraction means hiding the internal implementation details and exposing only what's necessary. In this example, the user interacts with \texttt{bookRide()} without needing to know how it works behind the scenes.
 \item Polymorphism
\end{itemize}
\vspace{0.3cm}
\noindent
\textbf{Question 8:} What does inheritance allow a class to do?
\\[0.2cm]
\begin{itemize}
 \item Hide internal data from other classes.
 \item Create multiple unrelated classes.
 \item \textbf{[CORRECT]} Use properties and methods of another class.
\\[0.1cm]
\textit{Explanation:} 
Inheritance allows a class (child/subclass) to acquire attributes and behaviors (methods) from another class (parent/superclass).
 \item Destroy objects when they're no longer needed.
\end{itemize}
\vspace{0.3cm}
\noindent
\textbf{Question 9:} Which real-world example best illustrates inheritance?
\\[0.2cm]
\begin{itemize}
 \item A company owning several branches.
 \item A car having four wheels.
 \item A button on a webpage responding to a click.
 \item \textbf{[CORRECT]} A dog being a type of animal and sharing common animal traits.
\\[0.1cm]
\textit{Explanation:} 
A dog inherits traits from the animal class---this is similar to how a subclass inherits from a superclass.
\end{itemize}
\vspace{0.3cm}
\noindent
\textbf{Question 10:} What best demonstrates polymorphism?
\\[0.2cm]
\begin{itemize}
 \item A teacher teaching math, science, and history.
 \item \textbf{[CORRECT]} A shape class with \texttt{draw()} method, and circle, square, triangle each having its own \texttt{draw()}.
\\[0.1cm]
\textit{Explanation:} 
Different shapes implement their own version of the same method, \texttt{draw()}, demonstrating polymorphism.
 \item A class hiding data from the user.
 \item A car with four wheels and each wheel using the same \texttt{steer()} method.
\end{itemize}



% End of section content